\origpage{338}{N. Robertson, P.D. Seymour / Journal of Combinatorial Theory, Series B 90 (2004) 325--385}
\textbf{Proof.} Suppose that for some $n$ there is no such $h$. Then $c(\Sigma_\chi) \ge 2$ and $dist(\Gamma_i) \le n$ for infinitely many values of $i$, and we may assume that $dist(\Gamma_i) \le n$ for all $i \ge 1$, by replacing our original sequence by the appropriate subsequence. For each $i \ge 1$, let $F_i$ be a path of $\Gamma_i^*$ with $\frac{1}{2}|E(F_i)| \le n$ joining distinct poles, such that every vertex of $F_i$ is in $\Sigma_\chi \setminus bd(\Sigma_\chi)$ except the first two and the last two. Then $K_i = F_i \cap \Sigma_\chi$ is a feature in $\Phi_\chi$ of size $\le 4n$, and by 5.2 (by replacing our sequence by a subsequence) we may assume that all the $K_i$'s are equivalent. By 5.1 we may assume that all the $K_i$'s are equal, to some $K$ say. \fix{Since $V(K)$ has only finitely many subsets, we may also assume (by replacing our sequence by a subsequence) that $V(K) \cap V(\Gamma_i)$ is the same for all $i \ge 1$.}\fixnote{fix:5.4-same-vertices}{This assumption is used in the details added at the end of the proof (note~\ref{fix:5.4-proof-added}).} Let $\Sigma_1$ be the (unique) fragment obtained by cutting $\Sigma_\chi$ along $K$, with associated surjection $\psi$. Let $\chi'$ be the colour scheme defined by $\Sigma_{\chi'} = \Sigma_1$, $\Phi_{\chi'} = \psi^{-1}(\Phi_\chi) \cap \Sigma_1$, $\Omega_{\chi'}(k) = \Omega_\chi(k)$ $(k = 2, 3)$, and for each long side $S'$ of $\Phi_{\chi'}$, $\Omega_{\chi'}(S') = \Omega_\chi(S)$ where $S$ is the long side of $\Phi_\chi$ with $\psi(S') \subseteq S$. Then $\chi'$ is not orientedly bad, since $\chi$ satisfies \textbf{S}$_3$. For each $i \ge 1$, let $\Gamma_i' = \psi^{-1}(\Gamma_i) \cap \Sigma_1$, and for each $e \in E(\Gamma_i')$, let $\phi_i'(e) = \phi_i(e)$ if $e \in E(\Gamma_i)$, and $\phi_i'(e) = e$ if $e \notin E(\Gamma_i)$ (so that $e$ is a short side of $\psi^{-1}(\Phi_\chi) \cap \Sigma_1$). Then $(\Gamma_i', \phi_i')$ is a $\chi'$-coloured painting, by 5.3, and the sequence $(\Gamma_i', \phi_i')$ $(i = 1, 2, \dots)$ is similarly oriented (for if $\Sigma_{\chi'}$ is orientable then so is $\Sigma_\chi$). Thus, since $\chi'$ is not orientedly bad, there exist $j > i \ge 1$ and a linear inflation $\sigma'$ of $(\Gamma_i', \phi_i')$ in $(\Gamma_j', \phi_j')$. Define $\sigma$ by
\begin{align*}
\sigma(e) &= \psi(\sigma'(\psi^{-1}(e))) \quad (e \in E(\Gamma_i)), \\
V(\sigma(v)) &= \bigcup_{v' \in \psi^{-1}(v)} \psi(V(\sigma'(v'))) \quad (v \in V(\Gamma_i)),
\end{align*}
where each $\sigma(v)$ is an induced subdrawing of $sk(\Gamma_j)$ with vertex set as given. If $v \in V(\Gamma_i)$ satisfies $|\psi^{-1}(v)| > 1$, then every $v' \in \psi^{-1}(v)$ belongs to $V(\Phi_{\chi'})$, and consequently satisfies $v' \in V(\sigma'(v'))$. So $v$ belongs to $\psi(V(\sigma'(v')))$ for each such $v'$, and it follows that $\sigma(v)$ is connected. It is easy to deduce that $\sigma$ is a linear inflation of $(\Gamma_i, \phi_i)$ in $(\Gamma_j, \phi_j)$ (for similar arguments, see for example Section 8 of [3]). \fix{We give the details. Let $\Phi' = \Phi_{\chi'}$. We use the following facts about $\psi$, which follow from the definitions; they hold for $\Gamma_j$ as well as for $\Gamma_i$. Each point of $\Sigma_\chi \setminus U(K)$ has exactly one point of $\Sigma_1$ over it, and $\psi^{-1}(U(K)) \subseteq bd(\Sigma_1)$. Each edge of $\Gamma_i'$ not included in $bd(\Sigma_1)$ is $\psi^{-1}(e)$ for an internal edge $e$ of $\Gamma_i$; it has the same size as $e$, and $\psi$ maps $\gamma_{\Gamma'_i}(\psi^{-1}(e))$ to $\gamma_{\Gamma_i}(e)$. If $e$ is a border edge of $\Gamma_i$ then $\psi^{-1}(e)$ is a border edge of $\Gamma_i'$. An edge of $\Gamma_i'$ included in $bd(\Sigma_1)$ that has a point over $bd(\Sigma_\chi)$ lies over $bd(\Sigma_\chi)$ and is $\psi^{-1}(e)$ for a border edge $e$ of $\Gamma_i$, $\psi$ is one-to-one on its closure, and $\psi$ maps its march to that of $e$, since both edges are directed as the sides that include them and $\psi$ maps the direction of each long side of $\Phi'$ onto that of a side of $\Phi_\chi$. An edge of $\Gamma_i'$ included in $bd(\Sigma_1)$ with no point over $bd(\Sigma_\chi)$ lies over $U(K)$, and both its ends are in $V(\Phi')$. Finally, since $\sigma'$ respects $\Phi'$, we have $y \in V(\sigma'(y))$ for every $y \in V(\Phi')$, as we saw in Section 3; and so, since the $\sigma'(y)$ are mutually disjoint, if $x \in V(\Phi')$ and $x \in V(\sigma'(y))$ for some $y \in V(\Gamma_i')$, then $x = y$.

Let $e \in E(\Gamma_i)$. If $e$ is internal, then $\psi^{-1}(e)$ is an internal edge of $\Gamma_i'$, and so is $\sigma'(\psi^{-1}(e))$, since $\sigma'$ is linear; hence $\sigma'(\psi^{-1}(e)) = \psi^{-1}(e')$ for an internal edge $e'$ of $\Gamma_j$, and $\sigma(e) = e'$. If $e$ borders a side $S$ of $\Phi_\chi$, then $\psi^{-1}(e)$ borders the side $R$ of $\Phi'$ with $\psi(R) \subseteq S$, and so does $\sigma'(\psi^{-1}(e))$, since $\sigma'$ respects $\Phi'$; hence $\sigma(e)$ is a border edge of $\Gamma_j$ bordering $S$, and $\sigma'(\psi^{-1}(e)) = \psi^{-1}(\sigma(e))$. In both cases $\sigma(e)$ has the same size as $e$, $\psi$ maps the $k$th end of $\sigma'(\psi^{-1}(e))$ to the $k$th end of $\sigma(e)$ for each $k$, and $\phi_i(e) \le \phi_j(\sigma(e))$, since $\phi_i'(\psi^{-1}(e)) \le \phi_j'(\sigma'(\psi^{-1}(e)))$. Since $\sigma'$ is one-to-one on edges, so is $\sigma$. Also $\sigma$ respects $\Phi_\chi$.

Next we show that $V(\sigma(v)) \subseteq V(\Gamma_j)$ for $v \in V(\Gamma_i)$. Let $y \in \psi^{-1}(v)$ and $x \in V(\sigma'(y))$. Then $x \in V(\Gamma_j')$, so $\psi(x) \in V(\Gamma_j) \cup V(K)$. If $\psi(x) \notin V(\Gamma_j)$, then $x \in V(\Phi')$, so $x = y$ and $\psi(x) = v \in V(\Gamma_i) \cap V(K)$; but $V(K) \cap V(\Gamma_i) = V(K) \cap V(\Gamma_j)$, and so $\psi(x) \in V(\Gamma_j)$, a contradiction. Now we show that $\psi(V(\sigma'(y)))$ induces a connected subdrawing of $sk(\Gamma_j)$, which the argument above uses. Let $a, b \in V(\sigma'(y))$ be ends of an edge $f$ of $\Gamma_j'$, with $\psi(a) \ne \psi(b)$. If $f$ lies over $U(K)$, then $a, b \in V(\Phi')$, and so $a = y = b$, which is impossible. Otherwise $\psi(f)$ is an edge of $\Gamma_j$ with ends $\psi(a), \psi(b)$. So $\psi$ maps each path of $\sigma'(y)$ to a walk of $sk(\Gamma_j)$ with vertices in $\psi(V(\sigma'(y)))$, as required.

For disjointness, let $v_1, v_2 \in V(\Gamma_i)$ be distinct, and suppose that $z \in V(\sigma(v_1)) \cap V(\sigma(v_2))$; so for $k = 1, 2$ there exist $y_k \in \psi^{-1}(v_k)$ and $x_k \in V(\sigma'(y_k))$ with $\psi(x_k) = z$. If $z \notin U(K)$ then $x_1 = x_2$, and so $y_1 = y_2$ and $v_1 = v_2$, a contradiction. If $z \in U(K)$, then $z \in V(K)$, since $z \in V(\Gamma_j)$ and $U(\Gamma_j) \cap U(K) = V(\Gamma_j) \cap V(K)$; so $x_1, x_2 \in V(\Phi')$, hence $x_k = y_k$ for $k = 1, 2$, and $v_1 = z = v_2$, a contradiction. For the ends, let $v$ be the $k$th end of $e \in E(\Gamma_i)$, and let $y$ be the $k$th end of $\psi^{-1}(e)$. Then $\psi(y) = v$, and $\sigma'(y)$ contains the $k$th end $x$ of $\sigma'(\psi^{-1}(e))$; so $\sigma(v)$ contains $\psi(x)$, which is the $k$th end of $\sigma(e)$. Thus $\sigma$ is an inflation of $(\Gamma_i, \phi_i)$ in $(\Gamma_j, \phi_j)$.

It remains to show that $\sigma$ is linear. For the first condition, internal edges go to internal edges, and an edge bordering a side $S$ goes to an edge bordering $S$, and so bordering the same cuff. For the third, let $e$ border a cuff $\Theta$, let $S$ be the side with $e \subseteq S$, and orient $\Theta$ so that the tail of $e$ immediately precedes $e$. Since $e$ is directed as $S$, this orientation traverses $S$ from its tail to its head; and since $\sigma(e) \subseteq S$ is also directed as $S$, the tail of $\sigma(e)$ immediately precedes $\sigma(e)$. For the second, let $e'$ be a border edge of $\Gamma_j$ with $e' \notin \sigma(E(\Gamma_i))$. Then $\psi^{-1}(e')$ is a border edge of $\Gamma_j'$. It is not $\sigma'(f_1)$ for any $f_1 \in E(\Gamma_i')$; for otherwise $f_1$ is a border edge, since $\sigma'$ is linear, and borders the same side $R$ of $\Phi'$ as $\psi^{-1}(e')$, since $\sigma'$ respects $\Phi'$; and $R$ lies over a side of $\Phi_\chi$, so $f_1 = \psi^{-1}(f)$ for a border edge $f$ of $\Gamma_i$, and $\sigma(f) = e'$, a contradiction. Since $\sigma'$ is linear, $\psi^{-1}(e')$ is an edge of $\sigma'(y)$ for some $y \in V(\Gamma_i')$. If $\psi(y) \notin V(\Gamma_i)$, then $y \in V(\Phi')$ and $\psi(y) \in V(K) \setminus V(\Gamma_i) = V(K) \setminus V(\Gamma_j)$; so every edge of $\Gamma_j'$ incident with $y$ lies over $U(K)$, and its other end is in $V(\Phi')$ and so not in $V(\sigma'(y))$, contrary to $\psi^{-1}(e')$ being an edge of $\sigma'(y)$. Hence $\psi(y) = v \in V(\Gamma_i)$, both ends of $e'$ belong to $V(\sigma(v))$, and $e'$ is an edge of $\sigma(v)$, since $\sigma(v)$ is an induced subdrawing. So $\sigma$ is linear.}\fixnote{fix:5.4-proof-added}{The paper leaves these details to the reader. The argument above is the one used in the formalization.} This is a contradiction. $\square$

Let $\Gamma$ be a 2-cell painting in $\Sigma$ with $E(\Gamma) \ne \emptyset$. If $\hat{\Sigma}$ is not a sphere, we define $rep(\Gamma)$ to be the minimum of $\frac{1}{2}|E(F)|$ over all non-null-homotopic circuits $F$ of $\Gamma^*$. (This exists, by Theorem 11.10 of [2].) If $\hat{\Sigma}$ is a sphere we set $rep(\Gamma) = \infty$.

\textbf{5.5.} \textit{Let $\chi$ satisfy} \textbf{S}$_1$, \textbf{S}$_3$, \textit{and let} $(\Gamma_i, \phi_i)$ $(i = 1, 2, \dots)$ \textit{be a similarly oriented bad sequence for $\chi$. Then for all $n \ge 0$ there exists $h \ge 1$ such that $rep(\Gamma_i) > n$ for all $i \ge h$.}

\textbf{Proof.} Suppose that for some $n$ there is no such $h$. Then $\hat{\Sigma}_\chi$ is not a sphere, and for infinitely many $i \ge 1$ there is a circuit $F_i$ of $\Gamma_i^*$ with $\frac{1}{2}|E(F_i)| \le n$ such that $U(F_i)$ bounds no disc in $\hat{\Sigma}_\chi$. If we choose $F_i$ with $|V(F_i) \cap \Sigma_\chi|$ minimum, it is easy to see that for each pole $r^*$, $V(F)$ contains at most two neighbours of $r^*$, and at most one if $r^* \notin V(F_i)$. \fix{To see this, let $F = F_i$, let $r^*$ be a pole, and let $\Theta = \bar{r} \setminus r$. Every neighbour of $r^*$ in $\Gamma_i^*$ is a vertex of $\Gamma_i$ in $\Theta$, and every neighbour in $\Gamma_i^*$ of a vertex of $\Gamma_i$ in $\Theta$, other than $r^*$, lies in $\Sigma_\chi$. Let $a, b \in V(F)$ be distinct, let $R$ be a path of $\Gamma_i^*$ between $a$ and $b$ with $E(R) \cap E(F) = \emptyset$ and $V(R) \cap V(F) = \{a, b\}$, and let $P, Q$ be the two paths of $F$ between $a$ and $b$. Then $U(P \cup R)$ or $U(Q \cup R)$ bounds no disc in $\hat{\Sigma}_\chi$, for otherwise $U(P \cup R)$ and $U(Q \cup R)$ are both null-homotopic, and so is $U(F)$, since a closed walk around $F$ is homotopic to a closed walk around $P \cup R$ followed by one around $R \cup Q$. Suppose first that $r^* \notin V(F)$ and $V(F)$ contains two distinct neighbours $u, v$ of $r^*$. Take $a = u$, $b = v$, and let $R$ be the path with edges $ur^*$ and $r^*v$; and we may assume that $U(P \cup R)$ bounds no disc. Now $|E(Q)| \ge 2 = |E(R)|$, since $\Gamma_i^*$ is bipartite and $u, v \in V(\Gamma_i)$, and so $\frac{1}{2}|E(P \cup R)| \le \frac{1}{2}|E(F)| \le n$. Moreover the neighbour of $u$ in $Q$ is not $r^*$, so it lies in $\Sigma_\chi$, and it is not in $P \cup R$; and $r^* \notin \Sigma_\chi$. Hence $|V(P \cup R) \cap \Sigma_\chi| < |V(F) \cap \Sigma_\chi|$, contrary to the choice of $F_i$. Now suppose that $r^* \in V(F)$, and that $v \in V(F)$ is a neighbour of $r^*$ in $\Gamma_i^*$ such that no edge of $F$ has ends $v, r^*$. Take $a = v$, $b = r^*$, and let $R$ consist of one edge of $\Gamma_i^*$ with ends $v, r^*$; and again we may assume that $U(P \cup R)$ bounds no disc. Now $|E(Q)| \ge 2 > |E(R)|$, since no edge of $F$ has ends $v, r^*$; and the neighbour of $v$ in $Q$ is not $r^*$, so it lies in $\Sigma_\chi$, and it is not in $P \cup R$. As before, this contradicts the choice of $F_i$. Hence every neighbour of $r^*$ in $V(F)$ is joined to $r^*$ by an edge of $F$, and there are at most two such, since $r^*$ has degree two in $F$.}\fixnote{fix:5.5-proof-added}{The paper leaves this argument to the reader. The argument above is the one used in the formalization.} As in 5.4 we may assume that $F_i$ exists for all $i \ge 1$, and all the $F_i$ are equal to some $F$ say. Let $F \cap \Sigma_\chi = K$. Then $K$ is a feature. \fix{As in 5.4, we may also assume that $V(K) \cap V(\Gamma_i)$ is the same for all $i \ge 1$.}\fixnote{fix:5.5-same-vertices}{This assumption is used in the details added at the end of the proof (note~\ref{fix:5.5-linear-proof-added}).}
